\documentclass[10pt,a4paper]{amsart}
\usepackage{amsmath,amssymb,amsthm,graphicx}
\usepackage{verbatim}
\newenvironment{reference}{\comment}{\endcomment}
%\newenvironment{reference}{}{}
\newenvironment{slogan}{\comment}{\endcomment}
\newenvironment{history}{\comment}{\endcomment}

% For commutative diagrams we use Xy-pic
\usepackage[all]{xy}

% We use 2cell for 2-commutative diagrams.
\xyoption{2cell}
% Theorem environments.
%
\theoremstyle{plain}
\newtheorem{theorem}[subsection]{Theorem}
\newtheorem{proposition}[subsection]{Proposition}
\newtheorem{lemma}[subsection]{Lemma}

\theoremstyle{definition}
\newtheorem{definition}[subsection]{Definition}
\newtheorem{example}[subsection]{Example}
\newtheorem{exercise}[subsection]{Exercise}
\newtheorem{situation}[subsection]{Situation}

\theoremstyle{remark}
\newtheorem{remark}[subsection]{Remark}
\newtheorem{remarks}[subsection]{Remarks}

\numberwithin{equation}{subsection}

% Macros
%
\def\lim{\mathop{\mathrm{lim}}\nolimits}
\def\colim{\mathop{\mathrm{colim}}\nolimits}
\def\Spec{\mathop{\mathrm{Spec}}}
\def\Hom{\mathop{\mathrm{Hom}}\nolimits}
\def\Ext{\mathop{\mathrm{Ext}}\nolimits}
\def\SheafHom{\mathop{\mathcal{H}\!\mathit{om}}\nolimits}
\def\SheafExt{\mathop{\mathcal{E}\!\mathit{xt}}\nolimits}
\def\Sch{\mathit{Sch}}
\def\Mor{\mathop{\mathrm{Mor}}\nolimits}
\def\Ob{\mathop{\mathrm{Ob}}\nolimits}
\def\Sh{\mathop{\mathit{Sh}}\nolimits}
\def\NL{\mathop{N\!L}\nolimits}
\def\CH{\mathop{\mathrm{CH}}\nolimits}
\def\proetale{{pro\text{-}\acute{e}tale}}
\def\etale{{\acute{e}tale}}
\def\QCoh{\mathit{QCoh}}
\def\Ker{\mathop{\mathrm{Ker}}}
\def\Im{\mathop{\mathrm{Im}}}
\def\Coker{\mathop{\mathrm{Coker}}}
\def\Coim{\mathop{\mathrm{Coim}}}

% Boxtimes
%
\DeclareMathSymbol{\boxtimes}{\mathbin}{AMSa}{"02}

%
% Macros for moduli stacks/spaces
%
\def\QCohstack{\mathcal{QC}\!\mathit{oh}}
\def\Cohstack{\mathcal{C}\!\mathit{oh}}
\def\Spacesstack{\mathcal{S}\!\mathit{paces}}
\def\Quotfunctor{\mathrm{Quot}}
\def\Hilbfunctor{\mathrm{Hilb}}
\def\Curvesstack{\mathcal{C}\!\mathit{urves}}
\def\Polarizedstack{\mathcal{P}\!\mathit{olarized}}
\def\Complexesstack{\mathcal{C}\!\mathit{omplexes}}
% \Pic is the operator that assigns to X its picard group, usage \Pic(X)
% \Picardstack_{X/B} denotes the Picard stack of X over B
% \Picardfunctor_{X/B} denotes the Picard functor of X over B
\def\Pic{\mathop{\mathrm{Pic}}\nolimits}
\def\Picardstack{\mathcal{P}\!\mathit{ic}}
\def\Picardfunctor{\mathrm{Pic}}
\def\Deformationcategory{\mathcal{D}\!\mathit{ef}}
\usepackage{cleveref}
\crefname{originaldefinition}{original definition}{original definitions}
\Crefname{originaldefinition}{Original definition}{Original definitions}
\makeatletter
\newlabel{definition-modules}{{07IS}{source}{}{sourcepointer-1}{}}
\newlabel{definition-modules@cref}{{[originaldefinition][0][]07IS}{[1][0][]source}{}{}{}}
\makeatother
\title{Quasi-coherent equals locally quasi-coherent crystal in O\_\{X/S\}-modules}
\author{The Stacks Project Authors}
\date{stacks-project:crystalline.tex:local-snapshot-20261012; GNU FDL 1.2 or later}
\begin{document}
\maketitle
\noindent\textit{Editorial note.} All delivered mathematical text is the author's own, copied verbatim from the pinned Stacks Project chapter file crystalline.tex; no statement or proof step was rewritten, reordered or removed.\par\medskip
\noindent\textit{Editorial note.} Span recomputation by label pairing (LANE\_BRIEF 10.2): in this chapter the author's proof environment follows the statement environment as a standalone block; the spec covers the complete environment pair (statement span, proof span) as recorded in the pool candidate.\par\medskip
\ifdefined\hypertarget\hypertarget{sourcepointer-1}{}\fi
\noindent\textit{Source reference.} Original definition 07IS: sources/stacks-chow/tags.txt, source lines 9597--9597. Cross-chapter standard prerequisite cited by the author inside the counted span; it is used as a pointer and its own proof is not part of this item.\par\medskip
\begin{lemma}
\label{lemma-crystal-quasi-coherent-modules}
With notation $X/S, \mathcal{I}, \gamma, \mathcal{C}, \mathcal{F}$
as in Definition \ref{definition-modules}. The following are equivalent
\begin{enumerate}
\item $\mathcal{F}$ is quasi-coherent, and
\item $\mathcal{F}$ is locally quasi-coherent and a crystal in
$\mathcal{O}_{X/S}$-modules.
\end{enumerate}
\end{lemma}

\begin{proof}
Assume (1). Let $f : (U', T', \delta') \to (U, T, \delta)$ be an object of
$\mathcal{C}$. We have to prove (a) $\mathcal{F}_T$ is a quasi-coherent
$\mathcal{O}_T$-module and (b) $c_f : f^*\mathcal{F}_T \to \mathcal{F}_{T'}$
is an isomorphism. The assumption means that we can find a covering
$\{(T_i, U_i, \delta_i) \to (T, U, \delta)\}$ and for each $i$
the restriction of $\mathcal{F}$ to $\mathcal{C}/(T_i, U_i, \delta_i)$
has a global presentation. Since it suffices to prove (a) and (b)
Zariski locally, we may replace $f : (T', U', \delta') \to (T, U, \delta)$
by the base change to $(T_i, U_i, \delta_i)$ and assume that $\mathcal{F}$
restricted to $\mathcal{C}/(T, U, \delta)$ has a global
presentation
$$
\bigoplus\nolimits_{j \in J}
\mathcal{O}_{X/S}|_{\mathcal{C}/(U, T, \delta)} \longrightarrow
\bigoplus\nolimits_{i \in I}
\mathcal{O}_{X/S}|_{\mathcal{C}/(U, T, \delta)} \longrightarrow
\mathcal{F}|_{\mathcal{C}/(U, T, \delta)}
\longrightarrow 0
$$
It is clear that this gives a presentation
$$
\bigoplus\nolimits_{j \in J} \mathcal{O}_T \longrightarrow
\bigoplus\nolimits_{i \in I} \mathcal{O}_T \longrightarrow
\mathcal{F}_T
\longrightarrow 0
$$
and hence (a) holds. Moreover, the presentation restricts to $T'$
to give a similar presentation of $\mathcal{F}_{T'}$, whence (b) holds.

\medskip\noindent
Assume (2). Let $(U, T, \delta)$ be an object of $\mathcal{C}$.
We have to find a covering of $(U, T, \delta)$ such that $\mathcal{F}$ has a
global presentation when we restrict to the localization of $\mathcal{C}$
at the members of the covering. Thus we may assume that $T$ is affine.
In this case we can choose a presentation
$$
\bigoplus\nolimits_{j \in J} \mathcal{O}_T \longrightarrow
\bigoplus\nolimits_{i \in I} \mathcal{O}_T \longrightarrow
\mathcal{F}_T
\longrightarrow 0
$$
as $\mathcal{F}_T$ is assumed to be a quasi-coherent $\mathcal{O}_T$-module.
Then by the crystal property of $\mathcal{F}$ we see that this pulls back
to a presentation of $\mathcal{F}_{T'}$ for any morphism
$f : (U', T', \delta') \to (U, T, \delta)$ of $\mathcal{C}$.
Thus the desired presentation of $\mathcal{F}|_{\mathcal{C}/(U, T, \delta)}$.
\end{proof}

\end{document}
